\section{Обмеження та загрози валідності}
\subsection{Внутрішня й конструктна валідність}
Оцінювання нерівномірне за задумом, бо нерівномірна реалізація. Точне відтворення графа шинглів вимірює чітко визначене обчислення, тоді як більшість семантичних завдань не мають незалежного еталона. Поточний збирач тверджень обумовлений обраними судженнями, запитами, джерелами й типами обсягу. Усунення прямого вставлення очікуваних міток є суттєвим покращенням, але не перетворює цей корпус на сліпий бенчмарк. Спеціально сконструйовані контрприклади встановлюють можливість збою, а не його частоту. Пізніший стрес-тест має те саме обмеження: його випадки написано після того, як збої стали відомими, тією самою ШІ-системою, яка переглянула правило. Це розрізнення стосується й зафіксованих синтетичних контролів звітів та подій.

Конструкти особливо вразливі до підміни. Лексичні теми не є позицією; лічильники реакцій не є схваленням кожного судження; повторений текст не є незалежним свідченням; інфраструктура, пов'язана з державою, не є операційною власністю; технічно дійсний кліп не є перевіреною сценою; а точні зміщення цитати не є випливанням. Надані корені ліній походження можуть запобігти відомому розгалуженню копій, але невідомі залежності лишаються. Захист фізичних фактів в агрегаторі перевірки залежить від правильних наданих міток обсягу, тож хибно класифікований обсяг усе ще може перебільшити висновок.

\begin{table}[!htbp]
\caption{Шари надійності та твердження, які лишаються невстановленими.}
\label{tab:reliability}\small
\begin{tabularx}{\linewidth}{@{}p{27mm}X X@{}}\toprule
Шар & Наявні докази & Що вони не встановлюють\\\midrule
Програмний & Успішний регресійний набір, схеми, бар'єри посилань і граничних часів & Повне покриття тестами, сертифікацію безпеки, семантичну правильність\\
Статистичний & Повторно обчислений точний граф, збережені вибірки, префікси із зерном, арифметика Голма & Специфічність щодо доброякісних подій, атрибуцію кампанії, причинний вплив\\
Доказовий & Зібрані байти, локатори, оголошені лінії походження, заміщення & Істинність чи автентичність джерел, незалежні корені, випливання з уривків\\
Мультимодальний & Квитанції декодування й обробки, порівняння тексту з наданим еталоном & Надійну особу, місце й час, незалежну безпеку цитування\\
Звітний & Хеші артефактів, перевірки цитат, посилань і чисел & Фактичну повноту, виявлення всіх непідтриманих тверджень\\\bottomrule
\end{tabularx}
\end{table}

\subsection{Зовнішня валідність і валідність статистичних висновків}
Центральні навантаження дає лише один реальний інцидент. Обрані канали, вибрані гілки, кураторські твердження й доступні медіа не утворюють репрезентативної рамки. Видалені матеріали, модерація, ранжування, пошуковий добір і обмеження платформ можуть впливати як на спостережений обсяг, так і на зміст. Жоден набір підібраних органічних подій не оцінює частоти хибнопозитивної координації. Наведені відхилення за рандомізацією залежать від обмінюваності в межах обраних страт; жодне не встановлює поширеності в генеральній сукупності чи причинних інформаційних ефектів. Спільні вузли й скопійовані джерела також унеможливлюють трактування ребер графа чи повторених тверджень як незалежних спостережень.

18 кліпів ASR обрані й нечисленні. Еталонні рядки, акустичні страти й позначки цитування не мають незалежно встановленого походження рішень. Різниці в частоті помилок між стратами можуть відображати лексичну складність і добір, а не лише шум. Перегляд семи активів більшою моделлю не є парним випробуванням продуктивності на всьому корпусі. Перцептивні пороги не мають розміченого бенчмарка збігів і незбігів. Негативний результат зіставлення медіа не може встановити, що записи зображують різні події.

\subsection{Валідність автоматизації, джерел і часу}
Оркестратор інтегрує обмежений адаптер збору й попередньо типізований аналіз, а не всі збирачі платформ. Фізичні згадки й частина вмісту звітів є кураторськими. Жоден незалежний наскрізний бенчмарк на рівні інцидентів не вимірює разом пошук джерел, виділення тверджень, достатність доказів, реконструкцію фізичних подій і правильність кінцевого звіту. Внутрішня перевірка ШІ-агентами не є незалежним людським рецензуванням; імена виконавців і мітки прийняття не встановлюють участі людини. Числові константи впевненості не є каліброваною невизначеністю.

Час збору й час оцінювання не можуть відновити, що вебсайт показував раніше. Збої HTTP/TLS залишають прогалини в охопленні. Деякі історичні сирі маніфести або приватні похідні відсутні у відстежуваних артефактах, що обмежує повторне виконання. Зібрана Source Watch копія, пов'язана зі звітом, та інша локально доступна копія мають лишатися окремими, навіть якщо кількості в заголовках збігаються. Час публікації, збору, запису, заявленої події й оцінювання не можна мовчки підміняти. Застаріла обробка часу для конкретного кейсу також обмежує повторне використання для інших сезонів і місць.

\subsection{Приватність, етика й межі оприлюднення}
Публічна доступність не усуває ризику для приватності. Псевдонімні ідентифікатори все ще можуть уможливлювати зв'язування, а дослівні коментарі, докладні хронології та поєднання між платформами можуть викривати окремих людей. Цей публікаційний пакет містить агреговані вимірювання, синтетичні випробування, посилання на код і хеші, а не сирі ідентифікатори коментаторів, номери телефонів, автентифікаційні матеріали чи масовий текст коментарів. Публічний репозиторій розробки йде далі: він містить звіти про кейс, які називають публічно ідентифікованих людей, і фрази, похідні від публічних коментарів. Таке оприлюднення є рішенням автора; автоматичне сканування не виявило облікових даних чи номерів телефонів, але незалежної перевірки приватності не проводилося. Поширення сирих даних потребує окремої перевірки приватності й умов платформ. Берклійський протокол і література про упередження соціальних даних підтримують розмежування етики отримання даних і якості доказів від простої технічної доступності \citep{berkeleyProtocol2022,Olteanu2019}.
